\section{Structure of a C++ Program}

In this section, we'll explore the fundamental building blocks of C++ programs. We'll start with a simple program and then break down the essential elements: keywords, identifiers, punctuation, operators, and syntax. Understanding these elements is essential for writing and reading C++ code effectively.

\subsection{Your First C++ Program}

Let's begin with the classic ``Hello, World!'' program---the traditional first program in any language:

\begin{lstlisting}
	#include <iostream>
	
	int main() {
		std::cout << "Hello, World!" << std::endl;
		return 0;
	}
\end{lstlisting}

Even this simple program contains several important elements. Let's walk through each line:

\begin{enumerate}
	\item \texttt{\#include <iostream>} — A \textbf{preprocessor directive} that includes the input/output stream library. This gives us access to \texttt{std::cout} for printing to the console.
	
	\item \texttt{int main() \{} — The \textbf{main function} declaration. Every C++ program must have exactly one \texttt{main} function—this is where program execution begins. The \texttt{int} indicates it returns an integer.
	
	\item \texttt{std::cout << "Hello, World!" << std::endl;} — This line outputs text to the console. \texttt{std::cout} is the standard output stream, \texttt{<<} is the stream insertion operator, and \texttt{std::endl} ends the line.
	
	\item \texttt{return 0;} --- Returns the value 0 to the operating system, indicating successful execution. By convention, returning 0 means ``no errors.''
	
	\item \texttt{\}} — Closes the main function's code block.
\end{enumerate}

\begin{remark}
	Don't worry if some of these concepts seem unfamiliar—we'll explore each one in detail throughout this section.
\end{remark}

\subsection{Overview of C++ Program Elements}

Before diving deeper into specific topics, let's understand the basic building blocks that make up every C++ program.

\subsubsection{Keywords}

Keywords are part of the vocabulary of a programming language. In most languages, keywords are also \textbf{reserved}, meaning you cannot redefine their meaning or use them in ways they were not intended.

C++ has approximately \textbf{90 keywords}, which is significantly more than many other languages:
\begin{itemize}
	\item C++: \textasciitilde{}90 keywords
	\item Java: \textasciitilde{}50 keywords
	\item C\@: 32 keywords
	\item Python: 33 keywords
\end{itemize}

In general, the more keywords a language has, the more complex its grammar. However, many C++ keywords are rarely used in everyday programming. You don't need to memorize all the keywords—you'll learn them naturally as you progress.

\begin{example}[Common C++ Keywords]{\leavevmode}
	Here are some frequently used keywords you'll encounter:
	\begin{itemize}
		\item \texttt{int}, \texttt{double}, \texttt{char}, \texttt{bool} — data type declarations
		\item \texttt{return} — returns a value from a function
		\item \texttt{if}, \texttt{else} — conditional statements
		\item \texttt{for}, \texttt{while}, \texttt{do} — loop constructs
		\item \texttt{class}, \texttt{struct} — defines user-defined types
		\item \texttt{const} — declares constants
		\item \texttt{public}, \texttt{private}, \texttt{protected} — access specifiers
		\item \texttt{virtual}, \texttt{override} — polymorphism support
		\item \texttt{new}, \texttt{delete} — dynamic memory management
		\item \texttt{try}, \texttt{catch}, \texttt{throw} — exception handling
	\end{itemize}
\end{example}

A complete list of C++ keywords can be found at \href{https://en.cppreference.com/w/cpp/keyword}{cppreference.com}.

\subsubsection{Identifiers}

While keywords are reserved by the language, \textbf{identifiers} are names that \textit{you}, the programmer, create. Identifiers represent meaningful entities in your program such as variables, functions, classes, and namespaces.

\paragraph{Identifier Naming Rules}

C++ has strict rules for what constitutes a valid identifier:

\begin{enumerate}
	\item \textbf{Must begin with a letter (a-z, A-Z) or underscore (\_)}
	\begin{itemize}
		\item Valid: \texttt{count}, \texttt{\_value}, \texttt{MyClass}
		\item Invalid: \texttt{2ndPlace}, \texttt{@name}
	\end{itemize}
	
	\item \textbf{Can contain letters, digits (0--9), and underscores}
	\begin{itemize}
		\item Valid: \texttt{player1}, \texttt{total\_score}, \texttt{calculateArea2D}
		\item Invalid: \texttt{my-variable}, \texttt{user name}
	\end{itemize}
	
	\item \textbf{Cannot be a C++ keyword}
	\begin{itemize}
		\item Invalid: \texttt{int}, \texttt{class}, \texttt{return}
	\end{itemize}
	
	\item \textbf{Case-sensitive}
	\begin{itemize}
		\item \texttt{myVariable}, \texttt{MyVariable}, and \texttt{MYVARIABLE} are three different identifiers
	\end{itemize}
	
	\item \textbf{No length limit} (though very long names are impractical)
\end{enumerate}

\begin{remark}
	Identifiers beginning with an underscore followed by an uppercase letter (e.g., \texttt{\_Max}) or containing double underscores (e.g., \texttt{my\_\_var}) are reserved for the compiler and standard library. Avoid using these patterns in your own code.
\end{remark}

\paragraph{Naming Conventions}

While not enforced by the compiler, following consistent naming conventions makes code more readable:

\begin{lstlisting}
	// Common conventions
	int playerScore;          // camelCase - common for variables
	int player_score;         // snake_case - also common for variables
	void CalculateTotal();    // PascalCase - common for functions/classes
	const int MAX_SIZE = 100; // SCREAMING_SNAKE_CASE - common for constants
\end{lstlisting}

Choose a convention and use it consistently throughout your project.

\subsubsection{Operators}

C++ has a rich set of operators beyond the standard arithmetic operations (addition, subtraction, multiplication, division). Let's look at some operators you'll encounter early on.

\begin{example}[Stream Operators]{\leavevmode}
	
	\begin{lstlisting}
		std::cout << "Hello, World!";  // << is the stream insertion operator
		std::cin >> favorite_number;   // >> is the stream extraction operator
	\end{lstlisting}
	
	\begin{itemize}
		\item \texttt{<<} (stream insertion operator) — inserts data into an output stream
		\item \texttt{>>} (stream extraction operator) — extracts data from an input stream
	\end{itemize}
\end{example}

\begin{remark}
	The \texttt{<<} and \texttt{>>} operators are actually the \textbf{bitwise left shift} and \textbf{bitwise right shift} operators in their original form. C++ \textit{overloads} these operators for use with streams—a concept we'll explore in depth in the Operator Overloading section. This is a powerful feature that lets you define how operators work with your own types.
\end{remark}

\begin{example}[Scope Resolution Operator]{\leavevmode}
	
	\begin{lstlisting}
		std::cout  // :: is the scope resolution operator
	\end{lstlisting}
	
	The \texttt{::} operator is used to access members of a namespace or class. Here, it accesses \texttt{cout} from the \texttt{std} namespace.
\end{example}

We'll explore operators in depth in the Statements and Operators section, where you'll learn about arithmetic, relational, logical, and many other types of operators.

\subsubsection{Punctuation}

C++ uses various punctuation marks to structure code:

\begin{itemize}
	\item \textbf{Semicolons (;)} — terminate statements
	\item \textbf{Curly braces (\{\})} — define code blocks
	\item \textbf{Parentheses (())} — group expressions and function parameters
	\item \textbf{Quotation marks} (\texttt{\textquotedbl\textquotedbl} and \texttt{\textquotesingle\textquotesingle}) --- define strings and characters
	\item \textbf{Angle brackets (<>)} — used with templates and include directives
	\item \textbf{Square brackets ([])} — array subscript and lambda captures
	\item \textbf{Comma (,)} — separates items in lists
	\item \textbf{Colon (:)} — used in inheritance, initializer lists, and labels
\end{itemize}

\begin{example}[Punctuation in Action]{\leavevmode}
	
	\begin{lstlisting}
		int main() {                           // Parentheses and opening brace
			std::cout << "Hello";              // Quotes, semicolon
			return 0;                          // Semicolon
		}                                      // Closing brace
	\end{lstlisting}
\end{example}

\subsubsection{Syntax}

When you combine keywords, identifiers, operators, and punctuation according to the rules of the language, you create \textbf{syntax}. Syntax refers to the structure and meaning of your code as understood by the compiler.

The compiler translates your code into machine code based on \textit{exactly} what you write—it doesn't guess or interpret your intentions. This may sound intimidating at first, but with practice, C++ syntax becomes natural.

\begin{example}[Valid vs. Invalid Syntax]{\leavevmode}
	
	\begin{lstlisting}
		// Valid syntax
		int x = 5;
		std::cout << x << std::endl;
		
		// Invalid syntax - missing semicolon
		int y = 10
		std::cout << y << std::endl;  // Compiler error!
	\end{lstlisting}
\end{example}

\begin{remark}
	Don't worry if syntax feels overwhelming at first. As you write more code and see more examples, the patterns will become familiar. The compiler is your friend—it will tell you when something is wrong, helping you learn the correct syntax along the way.
\end{remark}

\subsection{The \texttt{\#include} Preprocessor Directive}

The \texttt{\#include} directive is one of the first things you'll see in almost every C++ program. It's a \textbf{preprocessor directive}—an instruction that runs before the actual compilation begins.

\subsubsection{What Does \texttt{\#include} Do?}

When the preprocessor encounters \texttt{\#include}, it literally copies and pastes the entire contents of the specified file into your source code at that location. This allows you to use code defined elsewhere, such as the standard library.

\begin{lstlisting}
	#include <iostream>  // Include the I/O stream library
	#include <string>    // Include the string library
	#include <vector>    // Include the vector library
\end{lstlisting}

\subsubsection{Angle Brackets vs. Quotes}

There are two forms of \texttt{\#include}:

\begin{lstlisting}
	#include <iostream>      // Angle brackets: system/standard library headers
	#include "myheader.h"    // Quotes: your own header files
\end{lstlisting}

\begin{itemize}
	\item \textbf{Angle brackets} \texttt{< >}: The preprocessor searches in system directories where standard library headers are stored
	\item \textbf{Double quotes} \texttt{" "}: The preprocessor first searches in the current directory (where your source file is), then falls back to system directories
\end{itemize}

\begin{example}[Common Standard Library Headers]{\leavevmode}
	
	\begin{lstlisting}
		#include <iostream>   // Input/output streams (cout, cin)
		#include <string>     // String class
		#include <vector>     // Dynamic array container
		#include <cmath>      // Math functions (sqrt, pow, sin, cos)
		#include <fstream>    // File input/output
		#include <algorithm>  // Algorithms (sort, find, etc.)
	\end{lstlisting}
\end{example}

\begin{remark}
	Always include only the headers you need. Including unnecessary headers can slow down compilation and increase the size of your program.
\end{remark}

\subsection{Comments}

Comments are programmer-readable explanations in the source code—notes, annotations, or anything that adds meaning to what the program is doing.

\subsubsection{Why Comments Matter}

One crucial fact to understand: \textbf{comments never make it to the compiler}. The preprocessor strips out all comments before compilation, replacing each with a single space.

This may sound unusual—why write something the compiler never sees? That's precisely the point. Comments are for humans: for the next programmer who works on the code, or for yourself when you need to modify it later. Comments explain what you did and why you did it.

\subsubsection{C++ Comment Styles}

C++ has two styles of comments:

\paragraph{Single-Line Comments}

Begin with \texttt{//} and continue to the end of the line:

\begin{lstlisting}
	// This is a single-line comment
	
	int favorite_number;  // This stores the user's favorite number
	
	std::cout << "Enter a number: ";  // Prompt the user
\end{lstlisting}

\paragraph{Multi-Line Comments}

Begin with \texttt{/*} and end with \texttt{*/}:

\begin{lstlisting}
	/* This is a multi-line comment
	that spans across multiple lines.
	Everything between the markers is ignored. */
	
	/*
	* Author: Jude Eschete
	* Date: January 2026
	* Description: Program to calculate factorial
	*/
\end{lstlisting}

\subsubsection{Best Practices for Comments}

\paragraph{Write Self-Documenting Code}

The ideal code should be largely self-documenting. Use meaningful variable names, clear function names, and logical structure. Comments should supplement—not replace—clear code.

\paragraph{Don't Comment the Obvious}

Avoid comments that simply restate what the code already clearly expresses:

\begin{lstlisting}
	// BAD: Commenting the obvious
	return 0;  // return zero
	int sum = a + b;  // add a and b
	
	// GOOD: Explaining complex or non-obvious code
	// Using binary search for O(log n) performance on sorted data
	int index = binarySearch(sortedArray, target);
\end{lstlisting}

\paragraph{Explain the ``Why'', Not Just the ``What''}

Comments are most valuable when they explain \textit{why} code exists or \textit{why} a particular approach was chosen:

\begin{lstlisting}
	// We use a hash map instead of a tree map because
	// our access pattern is predominantly random lookups
	// with very few iterations over the entire collection
	std::unordered_map<int, std::string> cache;
\end{lstlisting}

\paragraph{Keep Comments Updated}

When you modify code, update the corresponding comments. Outdated comments are worse than no comments at all—they mislead future programmers.

\subsection{The \texttt{main()} Function}

Every C++ program must contain exactly one function named \texttt{main}. This function serves as the \textbf{entry point}—when you run a C++ program, the operating system calls \texttt{main()} to begin execution.

\subsubsection{Basic Form}

The simplest form of \texttt{main} is:

\begin{lstlisting}
	int main() {
		// Your code here
		return 0;
	}
\end{lstlisting}

The \texttt{int} before \texttt{main} indicates that the function returns an integer value to the operating system:
\begin{itemize}
	\item \texttt{return 0;} — Indicates successful execution
	\item \texttt{return 1;} (or any non-zero value) — Indicates an error occurred
\end{itemize}

The operating system can use this return value to determine whether the program completed successfully. This is especially useful in scripts and automated processes.

\subsubsection{Command-Line Arguments}

Programs often need to accept input from the command line. C++ provides an alternative form of \texttt{main} for this:

\begin{lstlisting}
	int main(int argc, char* argv[]) {
		// argc: argument count (number of arguments)
		// argv: argument vector (array of argument strings)
		return 0;
	}
\end{lstlisting}

\begin{example}[Using Command-Line Arguments]{\leavevmode}
	
	\begin{lstlisting}
		#include <iostream>
		
		int main(int argc, char* argv[]) {
			std::cout << "Number of arguments: " << argc << std::endl;
			
			for (int i = 0; i < argc; i++) {
				std::cout << "Argument " << i << ": " << argv[i] << std::endl;
			}
			
			return 0;
		}
	\end{lstlisting}
	
	If you compile this program as \texttt{myprogram} and run it with:
	\begin{lstlisting}[language=bash,numbers=none]
		./myprogram hello world       % Unix/macOS
		.\myprogram.exe hello world   % Windows
	\end{lstlisting}

	The output would be:
	\begin{lstlisting}[numbers=none]
		Number of arguments: 3
		Argument 0: ./myprogram
		Argument 1: hello
		Argument 2: world
	\end{lstlisting}

	Note that \texttt{argv[0]} is always the program name itself.
\end{example}

\subsubsection{What Happens Before and After \texttt{main()}}

While \texttt{main()} is where your code begins, C++ actually does some work before and after:

\begin{itemize}
	\item \textbf{Before \texttt{main()}}: Global objects are constructed, static variables are initialized
	\item \textbf{After \texttt{main()}}: Global objects are destroyed in reverse order of construction
\end{itemize}

This is why the order of global variable definitions can sometimes matter.

\subsection{Namespaces}

As programs grow larger and incorporate external libraries, the chance of naming conflicts increases. What if two libraries both define a function called \texttt{print()}? C++ solves this problem with \textbf{namespaces}.

\subsubsection{What is a Namespace?}

A namespace is a declarative region that provides a scope for identifiers inside it. Namespaces prevent name collisions by grouping related code under a unique name.

\begin{example}[The \texttt{std} Namespace]{\leavevmode}
	
	The C++ Standard Library places all its identifiers inside the \texttt{std} (standard) namespace:
	
	\begin{lstlisting}
		#include <iostream>
		#include <string>
		#include <vector>
		
		int main() {
			std::string name = "Alice";        // std::string
			std::vector<int> numbers;          // std::vector
			std::cout << name << std::endl;    // std::cout, std::endl
			return 0;
		}
	\end{lstlisting}
	
	Without namespaces, \texttt{string}, \texttt{vector}, \texttt{cout}, and thousands of other standard library identifiers would pollute the global namespace.
\end{example}

\subsubsection{Accessing Namespace Members}

There are three ways to access identifiers within a namespace:

\paragraph{1. Fully Qualified Names (Recommended)}

Use the scope resolution operator \texttt{::} each time:

\begin{lstlisting}
	std::cout << "Hello" << std::endl;
	std::string name = "World";
\end{lstlisting}

This is the most explicit and safest approach—it's always clear where each identifier comes from.

\paragraph{2. Using Declarations}

Import specific identifiers into the current scope:

\begin{lstlisting}
	using std::cout;
	using std::endl;
	
	cout << "Hello" << endl;  // Now we can use cout and endl directly
	std::string name;         // Other std members still need qualification
\end{lstlisting}

This is a good middle ground—you import only what you need.

\paragraph{3. Using Directives}

Import an entire namespace:

\begin{lstlisting}
	using namespace std;
	
	cout << "Hello" << endl;
	string name = "World";
	vector<int> numbers;
\end{lstlisting}

\begin{remark}
	While \texttt{using namespace std;} is convenient for small programs and learning, it's generally \textbf{discouraged in production code} and \textbf{should never appear in header files}. It defeats the purpose of namespaces by importing everything into the global scope, potentially causing name collisions.
\end{remark}

\subsubsection{Creating Your Own Namespaces}

You can create your own namespaces to organize your code:

\begin{lstlisting}
	namespace MyMath {
		const double PI = 3.14159265359;
		
		double square(double x) {
			return x * x;
		}
		
		double cube(double x) {
			return x * x * x;
		}
	}
	
	int main() {
		double area = MyMath::PI * MyMath::square(5.0);
		return 0;
	}
\end{lstlisting}

\subsection{Basic Input and Output (I/O) Using \texttt{cin} and \texttt{cout}}

Input and output are fundamental to most programs. C++ provides the \texttt{iostream} library for console I/O, with \texttt{cout} for output and \texttt{cin} for input.

\subsubsection{Output with \texttt{cout}}

\texttt{std::cout} (character output) is the standard output stream, typically connected to your console/terminal. Use the \textbf{stream insertion operator} \texttt{<<} to send data to it.

\begin{lstlisting}
	#include <iostream>
	
	int main() {
		std::cout << "Hello, World!";  // Output a string
		return 0;
	}
\end{lstlisting}

\paragraph{Chaining Output}

You can chain multiple \texttt{<<} operators together:

\begin{lstlisting}
	std::cout << "The answer is " << 42 << " and pi is " << 3.14159;
	// Output: The answer is 42 and pi is 3.14159
\end{lstlisting}

\paragraph{Outputting Different Data Types}

\texttt{cout} automatically handles different data types:

\begin{lstlisting}
	#include <iostream>
	#include <string>   // Required for std::string

	int age = 25;
	double salary = 50000.50;
	char grade = 'A';
	std::string name = "Alice";

	std::cout << "Name: " << name << std::endl;
	std::cout << "Age: " << age << std::endl;
	std::cout << "Salary: $" << salary << std::endl;
	std::cout << "Grade: " << grade << std::endl;
\end{lstlisting}

\paragraph{New Lines}

There are two ways to create new lines:

\begin{lstlisting}
	// Using std::endl (also flushes the buffer)
	std::cout << "Line 1" << std::endl;
	std::cout << "Line 2" << std::endl;
	
	// Using the newline character \n (faster, doesn't flush)
	std::cout << "Line 1\n";
	std::cout << "Line 2\n";
	
	// You can combine them
	std::cout << "Line 1\n" << "Line 2\n" << "Line 3" << std::endl;
\end{lstlisting}

\begin{remark}
	\texttt{std::endl} not only inserts a newline but also \textbf{flushes the output buffer}, ensuring all output is immediately written. This can be slower than \texttt{\textbackslash{}n}. Use \texttt{std::endl} when you need guaranteed immediate output (like before waiting for input), and \texttt{\textbackslash{}n} for performance-critical code.
\end{remark}

\begin{remark}
	The \texttt{<iostream>} library also provides \texttt{std::cerr} for error output and \texttt{std::clog} for logging. \texttt{std::cerr} is unbuffered, so error messages are always written immediately—useful when a crash might prevent a buffered stream from flushing.

	\begin{lstlisting}
		std::cout << "Normal output\n";
		std::cerr << "Error: something went wrong\n";
	\end{lstlisting}
\end{remark}

\subsubsection{Input with \texttt{cin}}

\texttt{std::cin} (character input) is the standard input stream, typically connected to your keyboard. Use the \textbf{stream extraction operator} \texttt{>>} to read data from it.

\begin{lstlisting}
	#include <iostream>
	
	int main() {
		int age;
		
		std::cout << "Enter your age: ";
		std::cin >> age;
		
		std::cout << "You are " << age << " years old." << std::endl;
		
		return 0;
	}
\end{lstlisting}

\paragraph{Reading Different Data Types}

\texttt{cin} automatically converts input to the appropriate type:

\begin{lstlisting}
	int whole_number;
	double decimal_number;
	char single_character;
	std::string word;
	
	std::cout << "Enter an integer: ";
	std::cin >> whole_number;
	
	std::cout << "Enter a decimal: ";
	std::cin >> decimal_number;
	
	std::cout << "Enter a character: ";
	std::cin >> single_character;
	
	std::cout << "Enter a word: ";
	std::cin >> word;  // Reads until whitespace
\end{lstlisting}

\paragraph{Chaining Input}

You can read multiple values in one statement:

\begin{lstlisting}
	int x, y, z;
	
	std::cout << "Enter three integers separated by spaces: ";
	std::cin >> x >> y >> z;
	
	std::cout << "You entered: " << x << ", " << y << ", " << z << std::endl;
\end{lstlisting}

The user can enter: \texttt{10 20 30} (separated by spaces or newlines).

\subsubsection{How \texttt{cin} Handles Input}

Understanding how \texttt{cin} works will save you from common bugs:

\begin{enumerate}
	\item \texttt{cin} skips leading whitespace (spaces, tabs, newlines)
	\item It reads until it encounters whitespace
	\item Remaining input stays in the \textbf{input buffer} for the next read
\end{enumerate}

\begin{example}[Input Buffer Behavior]{\leavevmode}
	
	\begin{lstlisting}
		std::string first, last;
		
		std::cout << "Enter your full name: ";
		std::cin >> first >> last;
		
		// If user enters "John Doe"
		// first = "John", last = "Doe"
	\end{lstlisting}
\end{example}

\subsubsection{Reading Entire Lines with \texttt{getline}}

Since \texttt{cin >>} stops at whitespace, it can't read phrases with spaces. Use \texttt{std::getline()} instead:

\begin{lstlisting}
	#include <iostream>
	#include <string>
	
	int main() {
		std::string full_name;
		
		std::cout << "Enter your full name: ";
		std::getline(std::cin, full_name);  // Reads the entire line
		
		std::cout << "Hello, " << full_name << "!" << std::endl;
		
		return 0;
	}
\end{lstlisting}

\paragraph{Mixing \texttt{cin >>} and \texttt{getline}}

A common pitfall occurs when mixing \texttt{cin >>} and \texttt{getline}:

\begin{lstlisting}
	int age;
	std::string name;
	
	std::cout << "Enter your age: ";
	std::cin >> age;  // Reads the number, leaves '\n' in buffer
	
	std::cout << "Enter your name: ";
	std::getline(std::cin, name);  // Reads the leftover '\n' immediately!
	// name is empty!
\end{lstlisting}

\textbf{Solution}: Clear the newline from the buffer:

\begin{lstlisting}
	std::cout << "Enter your age: ";
	std::cin >> age;
	std::cin.ignore(
		std::numeric_limits<std::streamsize>::max(), '\n');
	
	std::cout << "Enter your name: ";
	std::getline(std::cin, name);  // Now works correctly
\end{lstlisting}

Or use a simpler approach:

\begin{lstlisting}
	std::cin >> age;
	std::cin.ignore();  // Ignore one character (the newline)
	std::getline(std::cin, name);
\end{lstlisting}

\subsubsection{Input Validation}

\texttt{cin} can fail if the input doesn't match the expected type:

\begin{lstlisting}
	int number;
	
	std::cout << "Enter a number: ";
	std::cin >> number;
	
	if (std::cin.fail()) {
		std::cout << "Invalid input! That wasn't a number." << std::endl;
		std::cin.clear();  // Clear the error state
		std::cin.ignore(
			std::numeric_limits<std::streamsize>::max(), '\n');
	}
\end{lstlisting}

\begin{remark}
	The example below uses several concepts (\texttt{while} loops, conditionals) that are covered in later sections. Don't worry if it's not fully clear yet—we'll return to robust input handling once those tools are in your toolkit.
\end{remark}

\begin{example}[Robust Input Loop]{\leavevmode}

	\begin{lstlisting}
		#include <iostream>
		#include <limits>
		
		int main() {
			int age;
			
			while (true) {
				std::cout << "Enter your age: ";
				
				if (std::cin >> age && age > 0 && age < 150) {
					break;  // Valid input, exit loop
				}
				
				std::cout << "Invalid input. Please enter a valid age.\n";
				std::cin.clear();
				std::cin.ignore(
					std::numeric_limits<std::streamsize>::max(),
					'\n');
			}
			
			std::cout << "Your age is: " << age << std::endl;
			return 0;
		}
	\end{lstlisting}
\end{example}

\subsubsection{Common I/O Patterns}

Here are some patterns you'll use frequently:

\begin{lstlisting}
	#include <iostream>
	#include <string>

	int value;
	int x, y;
	std::string sentence;

	// Pattern 1: Prompt and read
	std::cout << "Enter value: ";
	std::cin >> value;

	// Pattern 2: Read multiple values
	std::cout << "Enter x and y: ";
	std::cin >> x >> y;

	// Pattern 3: Read a full line
	std::cout << "Enter a sentence: ";
	std::getline(std::cin, sentence);

	// Pattern 4: Echo back input
	std::cout << "You entered: " << value << std::endl;
\end{lstlisting}

\subsection{The C++ Preprocessor}

The C++ preprocessor is a program that processes your source code \textit{before} the compiler sees it. We've already seen \texttt{\#include}—now let's explore other preprocessor features.

\subsubsection{What Does the Preprocessor Do?}

The preprocessor performs several tasks:

\begin{enumerate}
	\item \textbf{Strips all comments} from the source file
	\item \textbf{Processes all \texttt{\#} directives}
	\item \textbf{Expands macros}
\end{enumerate}

\subsubsection{Common Preprocessor Directives}

\begin{itemize}
	\item \texttt{\#include} — includes the contents of another file
	\item \texttt{\#define} — defines a macro or constant
	\item \texttt{\#undef} — undefines a macro
	\item \texttt{\#ifdef} / \texttt{\#ifndef} — tests if a macro is defined
	\item \texttt{\#if} / \texttt{\#elif} / \texttt{\#else} / \texttt{\#endif} — conditional compilation
	\item \texttt{\#error} — generates a compilation error with a message
	\item \texttt{\#pragma} — implementation-specific directive
\end{itemize}

\subsubsection{The \texttt{\#define} Directive}

The \texttt{\#define} directive creates macros—text substitutions performed by the preprocessor.

\paragraph{Constant Macros}

\begin{lstlisting}
	#define PI 3.14159
	#define MAX_SIZE 100
	#define COMPANY_NAME "Acme Corp"
	
	double circumference = 2 * PI * radius;  // PI is replaced with 3.14159
	int buffer[MAX_SIZE];                    // MAX_SIZE is replaced with 100
\end{lstlisting}

\begin{remark}
	In modern C++, prefer \texttt{const} or \texttt{constexpr} variables over \texttt{\#define} for constants:
	
	\begin{lstlisting}
		// Modern C++ approach (preferred)
		constexpr double PI = 3.14159;
		constexpr int MAX_SIZE = 100;
	\end{lstlisting}
	
	These provide type safety and respect scope rules, unlike macros which are simple text replacement.
\end{remark}

\paragraph{Function-like Macros}

Macros can also take parameters:

\begin{lstlisting}
	#define SQUARE(x) ((x) * (x))
	#define MAX(a, b) ((a) > (b) ? (a) : (b))
	
	int result = SQUARE(5);      // Expands to: ((5) * (5))
	int larger = MAX(10, 20);    // Expands to: ((10) > (20) ? (10) : (20))
\end{lstlisting}

\begin{remark}
	Function-like macros have pitfalls—always wrap parameters in parentheses. In modern C++, prefer \texttt{inline} functions or \texttt{constexpr} functions instead.
\end{remark}

\subsubsection{Conditional Compilation}

Preprocessor directives can include or exclude code based on conditions:

\begin{lstlisting}
	#ifdef DEBUG
	std::cout << "Debug mode enabled" << std::endl;
	#endif
	
	#ifndef RELEASE
	// This code only compiles if RELEASE is not defined
	#endif
	
	#if defined(_WIN32)
	// Windows-specific code
	#elif defined(__APPLE__)
	// macOS-specific code
	#elif defined(__linux__)
	// Linux-specific code
	#else
	#error "Unsupported platform"
	#endif
\end{lstlisting}

\subsubsection{Include Guards}

When working with header files, include guards prevent multiple inclusion errors:

\begin{lstlisting}
	// myheader.h
	#ifndef MYHEADER_H
	#define MYHEADER_H
	
	// Header contents here
	class MyClass {
		// ...
	};
	
	#endif // MYHEADER_H
\end{lstlisting}

Or use the simpler, widely supported (but not officially standardized) alternative:

\begin{lstlisting}
	// myheader.h
	#pragma once
	
	// Header contents here
\end{lstlisting}

\subsubsection{The Preprocessor vs. The Compiler}

Here's a crucial point to understand:

\begin{center}
	\textbf{The C++ preprocessor does NOT understand C++!}
\end{center}

The preprocessor simply follows its own directives and prepares the source code. It doesn't parse C++ syntax, understand variables, or check types. The \textit{compiler} is the program that actually understands C++.

\subsection{The Compilation Process}

Understanding how C++ code is transformed into an executable helps you debug errors and understand the language better.

\subsubsection{From Source Code to Executable}

The build process consists of several stages:

\begin{enumerate}
	\item \textbf{Preprocessing}: Processes \texttt{\#} directives, removes comments, expands macros
	\item \textbf{Compilation}: Parses C++ code, checks syntax, generates assembly
	\item \textbf{Assembly}: Converts assembly to machine code (object files)
	\item \textbf{Linking}: Combines object files, resolves external references, creates executable
\end{enumerate}

\subsubsection{Compiling from the Command Line}

While IDEs handle compilation automatically, understanding command-line compilation is valuable:

\begin{lstlisting}[language=bash,numbers=none]
	# Compile a single file
	g++ -o myprogram main.cpp
	
	# Compile with warnings enabled (recommended)
	g++ -Wall -Wextra -o myprogram main.cpp
	
	# Compile with C++17 standard
	g++ -std=c++17 -Wall -o myprogram main.cpp
	
	# Compile multiple files
	g++ -std=c++17 -Wall -o myprogram main.cpp utils.cpp math.cpp
\end{lstlisting}

Common compiler flags:
\begin{itemize}
	\item \texttt{-o <filename>} — Specify output file name
	\item \texttt{-Wall} — Enable most warning messages
	\item \texttt{-Wextra} — Enable additional warnings
	\item \texttt{-std=c++17} — Use C++17 standard (or \texttt{c++20}, \texttt{c++23})
	\item \texttt{-g} — Include debugging information
	\item \texttt{-O2} — Enable optimizations
\end{itemize}

\subsubsection{Reading Compiler Error Messages}

When you make a mistake, the compiler will tell you. Learning to read error messages is one of the most valuable skills you can develop.

A typical GCC/Clang error looks like this:

\begin{lstlisting}[language=bash,numbers=none]
	main.cpp:5:5: error: expected ';' after return statement
		return 0
		        ^
		        ;
\end{lstlisting}

The format is: \texttt{filename:line:column: severity: message}. Here, line 5, column 5 of \texttt{main.cpp} is missing a semicolon. The caret (\texttt{\^{}}) points directly at the problem.

\begin{itemize}
	\item \textbf{error} — Must be fixed; compilation stops
	\item \textbf{warning} — Code compiles, but something is likely wrong; always investigate
	\item \textbf{note} — Additional context about an error above it
\end{itemize}

\begin{remark}
	When you have many errors, fix the \textit{first} one and recompile. A single mistake (like a missing semicolon or brace) can cause a cascade of confusing errors further down the file. After fixing the root cause, the others often disappear.
\end{remark}

\subsection{Header Files and Source Files}

As programs grow, organizing code across multiple files becomes essential. C++ uses a convention of separating code into \textbf{header files} and \textbf{source files}.

\subsubsection{The Separation of Declarations and Definitions}

\begin{itemize}
	\item \textbf{Declaration}: Tells the compiler that something exists and its type
	\item \textbf{Definition}: Provides the actual implementation
\end{itemize}

\begin{lstlisting}
	// Declaration (what it is)
	int add(int a, int b);
	
	// Definition (how it works)
	int add(int a, int b) {
		return a + b;
	}
\end{lstlisting}

\subsubsection{Header Files (\texttt{.h} or \texttt{.hpp})}

Header files contain declarations:

\begin{lstlisting}
	// math_utils.h
	#pragma once
	
	constexpr double PI = 3.14159265359;
	
	double circleArea(double radius);
	double circleCircumference(double radius);
\end{lstlisting}

\subsubsection{Source Files (\texttt{.cpp})}

Source files contain definitions:

\begin{lstlisting}
	// math_utils.cpp
	#include "math_utils.h"
	
	double circleArea(double radius) {
		return PI * radius * radius;
	}
	
	double circleCircumference(double radius) {
		return 2 * PI * radius;
	}
\end{lstlisting}

\subsubsection{Using Your Header Files}

\begin{lstlisting}
	// main.cpp
	#include <iostream>
	#include "math_utils.h"
	
	int main() {
		double radius = 5.0;
		std::cout << "Area: " << circleArea(radius) << std::endl;
		return 0;
	}
\end{lstlisting}

\subsubsection{Why Separate Headers and Source Files?}

\begin{enumerate}
	\item \textbf{Faster compilation}: Only recompile source files that changed
	\item \textbf{Code organization}: Clear separation of interface and implementation
	\item \textbf{Reduced dependencies}: Include only what you need
	\item \textbf{Information hiding}: Users see only the interface
\end{enumerate}

\subsection{Exercises}

\begin{example}[Exercise 1: Hello Message]{\leavevmode}
	
	Write a program that outputs ``Hello from C++!'' to the console.
	
	\textbf{Expected output:}
	\begin{lstlisting}[numbers=none]
		Hello from C++!
	\end{lstlisting}
\end{example}

\begin{example}[Exercise 2: Formatted Output]{\leavevmode}
	
	Write a program that outputs your name, age, and favorite programming language on separate lines, with labels.
	
	\textbf{Expected output (example):}
	\begin{lstlisting}[numbers=none]
		Name: Alice
		Age: 25
		Favorite Language: C++
	\end{lstlisting}
\end{example}

\begin{example}[Exercise 3: User Input]{\leavevmode}

	Write a program that:
	\begin{enumerate}
		\item Asks the user for their name
		\item Asks the user for their favorite number
		\item Outputs a personalized message with both pieces of information
	\end{enumerate}

	\textbf{Sample interaction:}
	\begin{lstlisting}[numbers=none]
		Enter your name: Bob
		Enter your favorite number: 42
		Hello, Bob! Your favorite number is 42.
	\end{lstlisting}
\end{example}

\begin{example}[Exercise 4: Namespaces]{\leavevmode}

	The following program uses \texttt{using namespace std;}. Rewrite it so that it:
	\begin{enumerate}
		\item Removes the \texttt{using namespace std;} directive
		\item Uses \texttt{using} declarations to import only \texttt{cout} and \texttt{endl}
		\item Uses a fully qualified name (\texttt{std::string}) for the string variable
	\end{enumerate}

	\textbf{Starting code:}
	\begin{lstlisting}
		#include <iostream>
		#include <string>
		using namespace std;

		int main() {
			string greeting = "Hello, C++!";
			cout << greeting << endl;
			return 0;
		}
	\end{lstlisting}
\end{example}